\item \model{MobileLLM-Pro}

\paragraph*{مقدمه و ساختار چندمرحله‌ای دادگان پیش‌آموزش}
مدل \model{MobileLLM-Pro} \cite{huber2025mobilellmpro} یک مدل زبانی بنیادین ۱ میلیارد پارامتری (دقیقاً ۱.۰۸۴ میلیارد پارامتر) است که با هدف استقرار و استنتاج کارآمد روی دستگاه‌های لبه‌ای (\en{On-Device}) طراحی شده است. از آنجا که مدل‌های با مقیاس زیر ۳ میلیارد پارامتر حساسیت بسیار بالایی نسبت به نویز دادگان و پدیده جابجایی توزیع (\en{Distribution Shift}) در طول مراحل گوناگون آموزش دارند، فرایند پیش‌آموزش داده‌ها در این مدل بر پایه یک خط‌لوله ۴ فازی مهندسی شده است:
\begin{itemize}
    \item \textbf{فاز ۱ (اکتساب زبان پایه - \en{Language Acquisition}):} تمرکز بر یادگیری عمومی زبان با بودجه آموزشی \textbf{۱.۴ تریلیون توکن} (\en{1.4T tokens}) و طول بافت $2\text{K}$ توکن.
    \item \textbf{فاز ۲ (بسط بافت متنی - \en{Context Expansion}):} توسعه پنجره بافت به \textbf{۱۲۸,۰۰۰ توکن} (\en{128K context}) روی \textbf{۲۰ میلیارد توکن} بدون تغییر در توزیع متنی فاز ۱.
    \item \textbf{فاز ۳ (بازپخت و ادغام متخصصان - \en{Specialist Model Merging}):} بازپخت مدل با خوشه‌بندی دادگان پیش‌آموزش روی حوزه‌های هدف با بودجه \textbf{۶۰ میلیون توکن} برای هر شاخه متخصص.
    \item \textbf{فاز ۴ (آموزش آگاه از کوانتایش - \en{Quantization-Aware Training}):} فشرده‌سازی و پایدارسازی وزن‌های مدل تا سطح ۴ بیتی بر روی \textbf{۸۰ میلیارد توکن} (حدود ۵٪ از بودجه کل آموزش).
\end{itemize}

در تمامی فازهای پیش‌آموزش، به‌جای استفاده از اهداف کلاسیک تک‌مقداری (\en{One-hot targets}) و تابع زیان انتروپی متقاطع مرسوم، از تقطیر دانش مبتنی بر لاجیت (\en{Logit-based Knowledge Distillation}) \cite{hinton2015distilling, jiao2020tinybert} از مدل معلم قدرتمند \en{Llama 4-Scout} با واژگان کامل ۲۰۲,۰۴۸ توکنی استفاده شده است. هدف بهینه‌سازی از طریق حداقل‌سازی واگرایی پیشرو کولبک-لایبلر (\en{Forward KL-Divergence}) تعریف می‌شود:
\begin{equation}
    \mathcal{L}_{\text{KD}}(\theta) = \mathcal{D}_{\text{KL}}\big(P_{\text{Teacher}}(y \mid x) \,\|\, P_{\text{Student}}(y \mid x; \theta)\big)
\end{equation}

\paragraph*{شبیه‌سازی آفلاین و گزینش بهینه داده‌ها (\en{SDM Data-Mix Simulations})}
به منظور پرهیز از تصمیم‌گیری‌های شهودی یا وزن‌دهی یکنواخت، ترکیب داده‌های فاز ۱ با استفاده از چهارچوب شبیه‌ساز مقیاس‌پذیر داده (\en{Scalable Data Mixer - SDM}) \cite{chang2025sdm, chang2025automixer} استخراج شده است. در این شیوه، اثرگذاری هر نمونه متنی $x$ بر دامنه‌های گوناگون داده‌ای $D$ (شامل زبان عمومی، دانش، استدلال، کد و ریاضیات) از طریق مدل‌های زبانی آماری سبک‌وزن (\en{SLM}) سنجیده می‌شود:
\begin{equation}
    \text{Inf}(x, D) = \mathcal{L}_D(\text{SLM}) - \mathcal{L}_D(\text{SLM} \leftarrow x)
\end{equation}
که در آن $\mathcal{L}_D$ زیان انتروپی متقاطع روی دامنه $D$ را قبل و بعد از آموزش یک‌گامه با نمونه $x$ نشان می‌دهد. این امتیازها در قالب یک بردار تأثیر $\mathbf{\Delta}_x = [\text{Inf}(x, D_1), \dots, \text{Inf}(x, D_k)]^T$ تجمیع شده و توسط یک رگرسور عصبی سبک‌وزن، مطلوبیت پایین‌دستی (\en{Downstream Utility}) $\hat{y}_D$ پیش‌بینی می‌شود. در نهایت، با یک بار پیمایش آفلاین کل پیکره در طول ۱ میلیون گام، وزن‌های ایستای نمونه‌برداری $w_i$ متناسب با مطلوبیت انتظاری $\mathbb{E}_{x \in \mathcal{D}_i}[\hat{y}]$ به دست می‌آیند به‌طوری که $\sum_i w_i = 1.0$.

آمار تفکیکی دادگان پالایش‌شده برای پیش‌آموزش فاز ۱ و ۲ در جدول زیر خلاصه شده است:

\begin{table}[htbp]
\centering
\caption{مشخصات آماری دادگان فاز ۱ و فاز ۲ پیش‌آموزش مدل \model{MobileLLM-Pro}}
\label{tab:mobilellm_pro_pretrain_data}
\resizebox{0.9\textwidth}{!}{%
\begin{tabular}{llccc}
\toprule
\textbf{مجموعه داده (\en{Dataset})} & \textbf{مرجع علمی} & \textbf{سطرها (میلیون)} & \textbf{حجم توکن (میلیارد)} & \textbf{وزن در آمیزه (\%)} \\
\midrule
\en{Fineweb-EDU (2024)} & Allal et al. \cite{allal2024fineweb} & 1279.1 & 1300.0 & \textbf{89.75\%} \\
\en{Starcoder (2023)} & Li et al. \cite{li2023starcoder} & 206.6 & 263.8 & \textbf{4.66\%} \\
\en{Open Web Math (2023)} & Paster et al. \cite{paster2023openwebmath} & 6.1 & 12.6 & \textbf{1.92\%} \\
\en{Arxiv (2023)} & - & 1.5 & 28.0 & \textbf{1.35\%} \\
\en{Wikipedia} & - & 7.2 & 3.7 & \textbf{1.02\%} \\
\en{Stack Exchange (2023)} & Together Computer \cite{together2023redpajama} & 29.2 & 19.6 & \textbf{1.02\%} \\
\en{Algebraic Stack (2023)} & Azerbayev et al. \cite{azerbayev2023llemma} & 3.4 & 12.6 & \textbf{0.24\%} \\
\midrule
\textbf{مجموع کل (\en{Total})} & - & \textbf{1500.0} & \textbf{1640.3} & \textbf{100.00\%} \\
\bottomrule
\end{tabular}%
}
\end{table}

این توزیع آماری نشان می‌دهد که مدل بیشترین سهم خود (۸۹.۷۵٪) را از متون آموزشی و منسجم وب دریافت کرده و تخصیص متوازن ۶.۸۲٪ از کل آمیزه به داده‌های استدلالی، برنامه‌نویسی و ریاضیات، ظرفیت محاسباتی مدل ۱ میلیارد پارامتری را بدون اشباع یا نویزپذیری پایدار ساخته است. این مرحله با اندازه دسته موثر ۲ میلیون توکن در ۶۴۰,۰۰۰ گام به‌روزرسانی و با نرخ یادگیری اولیه $4 \times 10^{-4}$ آموزش داده شده است.

\paragraph*{بسط بافت متن بدون جابجایی توزیع (\en{Implicit Positional Distillation})}
یکی از نوآوری‌های بنیادین در سازمان‌دهی داده‌های پیش‌آموزش \model{MobileLLM-Pro}، حذف نیاز به پیکره‌های متنی بلند مرسوم در فاز دوم است. تزریق متون فوق‌طولانی معمولاً تعادل زبانی مدل‌های فشرده را برهم می‌زند و به افت یادگیری منجر می‌شود. در این پژوهش، از تکنیک «تقطیر موقعیتی پنهان» بر بستر تعبیه موقعیتی چرخشی (\en{RoPE}) استفاده شده است:
\begin{equation}
    \begin{bmatrix} x'_i \\ x'_{i+1} \end{bmatrix} = \begin{bmatrix} \cos(\theta_i p) & -\sin(\theta_i p) \\ \sin(\theta_i p) & \cos(\theta_i p) \end{bmatrix} \begin{bmatrix} x_i \\ x_{i+1} \end{bmatrix}, \quad \theta_i = 10000^{-\frac{2i}{d}}
\end{equation}
با استفاده از اتصال اسناد کوتاه قبلی و پوشش علیتی بلوکی (\en{Block Causal Masking})، مدل دانشجو توزیع لاجیت معلم را دریافت می‌کند. از آنجا که مدل معلم قبلاً روی توالی‌های بلند آموزش دیده، لاجیت‌های آن رفتار چرخش زاویه‌ای کامل (تا سقف ۳۶۰ درجه معادل بافت ۱۲۸K) را به صورت پنهان بازتاب می‌دهند. در نتیجه، بدون مواجهه مستقیم مدل دانشجو با متون طولانی اختصاصی و تنها با آموزش روی \textbf{۲۰ میلیارد توکن مضاعف} از همان توزیع داده‌ای فاز اول (طی ۱۰۰,۰۰۰ گام با بیشینه نرخ یادگیری $4 \times 10^{-5}$)، مدل به پنجره بافت ۱۲۸K دست پیدا می‌کند.

\paragraph*{خوشه‌بندی داده‌ها و ادغام متخصصان (\en{Specialist Model Merging})}
در فاز سوم، داده‌های بدون ساختار پیش‌آموزش خوشه‌بندی شده و به زیرمجموعه‌های متمرکز بر حوزه‌های خاص تفکیک می‌شوند. به‌جای آموزش هم‌زمان که موجب رقابت مخرب بین قابلیت‌ها (نظیر حفظ حقایق در برابر توانایی استدلال) می‌شود، $n$ مدل متخصص موازی ($b_1, \dots, b_n$) هر کدام به‌طور مستقل روی دامنه متنی اختصاصی خود $\mathcal{D}_i$ به میزان \textbf{۶۰ میلیون توکن} طی \textbf{۵۰۰ گام} آموزش می‌بینند. در نهایت، پارامترهای متخصصان از طریق ادغام غیریکنواخت وزنی تلفیق می‌شوند:
\begin{equation}
    M = \frac{1}{n} \sum_{b=1}^{n} \theta_b \cdot w_b, \quad \sum_{b=1}^{n} w_b = 1.0
\end{equation}
این رویکرد داده‌محور، مدل پایه \model{MobileLLM-Pro-base} را حاصل می‌آورد که در آن ویژگی‌های مجزای حوزه‌ای به‌شکل هم‌افزا پیوند خورده‌اند.

\paragraph*{بودجه داده‌ها در فاز کوانتایش (\en{QAT Phase})}
در مرحله چهارم، به منظور انطباق نهایی وزن‌ها و فعال‌سازها برای استقرار روی پردازنده‌های کم‌مصرف \en{CPU} و شتاب‌دهنده‌های \en{NPU} (مانند \en{Apple ANE} و \en{Qualcomm HTP})، مدل با بودجه متنی معادل \textbf{۸۰ میلیارد توکن} با روش تقطیر از خود (\en{Self-Distillation}) از نسخه اعشاری فاز ۳، تحت یادگیری بازه کوانتایش قرار می‌گیرد.

\paragraph*{ارزیابی‌های تجربی و تحلیل آماری اثرگذاری داده‌ها}
جدول‌های زیر اثرات تجربی راهبردهای انتخاب، پالایش و عدم جابجایی توزیع داده‌ها را بر عملکرد نهایی مدل اثبات می‌کنند:

\begin{table}[htbp]
\centering
\caption{مقایسه عملکرد آمیزه داده بهینه‌شده با \en{SDM} در برابر آمیزه یکنواخت (\en{Uniform})}
\label{tab:mobilellm_pro_sdm_ablation}
\begin{tabular}{llcc}
\toprule
\textbf{مرحله آموزش} & \textbf{بودجه توکن} & \textbf{منبع آمیزه داده} & \textbf{میانگین عملکرد (\%) \en{(Avg. Score)}} \\
\midrule
\multirow{2}{*}{\textbf{پیش‌آموزش (\en{Pre-training})}} & \multirow{2}{*}{1.4T} & یکنواخت (\en{Uniform}) & 38.70\% \\
& & \textbf{\en{SDM} (پیشنهادی)} & \textbf{49.31\%} \\
\midrule
\multirow{2}{*}{\textbf{آموزش دستورپذیری (\en{IFT})}} & \multirow{2}{*}{80B} & یکنواخت (\en{Uniform}) & 17.94\% \\
& & \textbf{\en{SDM} (پیشنهادی)} & \textbf{45.23\%} \\
\bottomrule
\end{tabular}
\end{table}

بر اساس نتایج جدول فوق، بهینه‌سازی داده‌ها از طریق \en{SDM} موجب ارتقای مطلق \textbf{۱۰.۶۱٪} در مرحله پیش‌آموزش نسبت به توزیع یکنواخت شده است.

همچنین جدول زیر، مزیت ثابت نگه داشتن توزیع داده در فاز ۲ را در مقایسه با تزریق متون طولانی متداول اثبات می‌کند:

\begin{table}[htbp]
\centering
\caption{ارزیابی راهبردهای دادگان فاز ۲ در آزمون سوزن در انبار کاه (\en{NIH}) و میانگین تسک‌های عمومی}
\label{tab:mobilellm_pro_context_ablation}
\begin{tabular}{lcc}
\toprule
\textbf{پیکربندی داده‌ای فاز ۲} & \textbf{شاخص \en{NIH} (\%)} & \textbf{میانگین عملکرد عمومی (بدون \en{NIH})} \\
\midrule
فاز ۱ (طول بافت کوتاه $2\text{K}$) & 6.70\% & 53.74\% \\
فاز ۲ با متون طولانی متعارف & 80.22\% & 47.86\% (افت ۵.۸۸٪) \\
\textbf{فاز ۲ با تقطیر پنهان (توزیع داده‌های فاز ۱)} & \textbf{99.78\%} & \textbf{53.57\%} (حفظ ثبات کیفی) \\
\bottomrule
\end{tabular}
\end{table}

همان‌طور که مشاهده می‌شود، استفاده از دادگان متنی طولانی اختصاصی موجب افت ۵.۸۸٪ در عملکرد عمومی مدل گردیده، در حالی که روش تقطیر پنهان با حفظ توزیع داده‌های فاز ۱، نمره بازیابی \en{NIH} را به ۹۹.۷۸٪ رسانده و از افت عملکرد عمومی ممانعت نموده است.

در نهایت، جدول زیر نشان‌دهنده تأثیر خوشه‌بندی داده‌های پیش‌آموزش روی دامنه‌های تخصصی مستقل و سپس ادغام وزنی آن‌ها است:

\begin{table}[htbp]
\centering
\caption{تأثیر خوشه‌بندی دامنه‌ها روی متخصصان مجزا و هم‌افزایی حاصل از ادغام وزنی نهایی (\en{Weight Avg})}
\label{tab:mobilellm_pro_specialist_ablation}
\resizebox{\textwidth}{!}{%
\begin{tabular}{lcccccc}
\toprule
\textbf{بنچمارک (متریک)} & \textbf{پیش از بازپخت} & \textbf{متخصص ۱} & \textbf{متخصص ۲} & \textbf{متخصص ۳} & \textbf{متخصص ۴} & \textbf{ادغام نهایی (\en{Weight Avg.})} \\
\midrule
\en{HellaSwag (acc\_char)} & 65.77 & 66.69 & 66.87 & 66.31 & 67.04 & \textbf{67.58} \\
\en{BoolQ (acc\_char)} & 71.28 & 76.51 & 76.94 & 73.64 & 78.47 & \textbf{76.82} \\
\en{PIQA (acc\_char)} & 75.57 & 75.14 & 76.50 & 75.46 & 76.17 & \textbf{76.55} \\
\en{SIQA (acc\_char)} & 47.29 & 47.90 & 50.97 & 46.88 & 50.31 & \textbf{51.23} \\
\en{TriviaQA (em)} & 36.61 & 36.78 & 40.07 & 37.48 & 39.91 & \textbf{40.18} \\
\en{Natural Questions (em)} & 12.02 & 12.22 & 15.84 & 11.83 & 16.01 & \textbf{16.29} \\
\en{ARC-Challenge (acc\_char)} & 50.64 & 53.05 & 52.19 & 51.33 & 53.05 & \textbf{54.16} \\
\en{ARC-Easy (acc\_char)} & 71.37 & 75.22 & 76.45 & 73.66 & 76.74 & \textbf{76.96} \\
\en{WinoGrande (acc\_char)} & 63.30 & 63.06 & 63.14 & 62.90 & 63.06 & \textbf{63.54} \\
\en{OBQA (acc\_char)} & 41.80 & 41.40 & 43.40 & 41.80 & 43.40 & \textbf{44.00} \\
\en{NIH (em)} & 100.00 & 100.00 & 100.00 & 100.00 & 100.00 & \textbf{100.00} \\
\midrule
\textbf{میانگین (بدون \en{NIH})} & \textbf{53.57} & \textbf{54.80} & \textbf{56.24} & \textbf{54.13} & \textbf{56.42} & \textbf{56.73} \\
\bottomrule
\end{tabular}%
}
\end{table}

این آمار مؤید آن است که آموزش شاخه‌ای روی داده‌های تفکیک‌شده و ادغام متعاقب پارامترها، امتیاز مدل را از ۵۳.۵۷ به **۵۶.۷۳** ارتقا داده که حتی از بهترین متخصص منفرد (۵۶.۴۲) نیز فراتر رفته و هم‌افزایی بازنمایی دادگان را ثابت می‌کند.